\documentclass[10pt,a4paper]{amsart}
\usepackage{amsmath,amssymb,amsthm,graphicx}
\usepackage{verbatim}
\newenvironment{reference}{\comment}{\endcomment}
%\newenvironment{reference}{}{}
\newenvironment{slogan}{\comment}{\endcomment}
\newenvironment{history}{\comment}{\endcomment}

% For commutative diagrams we use Xy-pic
\usepackage[all]{xy}

% We use 2cell for 2-commutative diagrams.
\xyoption{2cell}
% Theorem environments.
%
\theoremstyle{plain}
\newtheorem{theorem}[subsection]{Theorem}
\newtheorem{proposition}[subsection]{Proposition}
\newtheorem{lemma}[subsection]{Lemma}

\theoremstyle{definition}
\newtheorem{definition}[subsection]{Definition}
\newtheorem{example}[subsection]{Example}
\newtheorem{exercise}[subsection]{Exercise}
\newtheorem{situation}[subsection]{Situation}

\theoremstyle{remark}
\newtheorem{remark}[subsection]{Remark}
\newtheorem{remarks}[subsection]{Remarks}

\numberwithin{equation}{subsection}

% Macros
%
\def\lim{\mathop{\mathrm{lim}}\nolimits}
\def\colim{\mathop{\mathrm{colim}}\nolimits}
\def\Spec{\mathop{\mathrm{Spec}}}
\def\Hom{\mathop{\mathrm{Hom}}\nolimits}
\def\Ext{\mathop{\mathrm{Ext}}\nolimits}
\def\SheafHom{\mathop{\mathcal{H}\!\mathit{om}}\nolimits}
\def\SheafExt{\mathop{\mathcal{E}\!\mathit{xt}}\nolimits}
\def\Sch{\mathit{Sch}}
\def\Mor{\mathop{\mathrm{Mor}}\nolimits}
\def\Ob{\mathop{\mathrm{Ob}}\nolimits}
\def\Sh{\mathop{\mathit{Sh}}\nolimits}
\def\NL{\mathop{N\!L}\nolimits}
\def\CH{\mathop{\mathrm{CH}}\nolimits}
\def\proetale{{pro\text{-}\acute{e}tale}}
\def\etale{{\acute{e}tale}}
\def\QCoh{\mathit{QCoh}}
\def\Ker{\mathop{\mathrm{Ker}}}
\def\Im{\mathop{\mathrm{Im}}}
\def\Coker{\mathop{\mathrm{Coker}}}
\def\Coim{\mathop{\mathrm{Coim}}}

% Boxtimes
%
\DeclareMathSymbol{\boxtimes}{\mathbin}{AMSa}{"02}

%
% Macros for moduli stacks/spaces
%
\def\QCohstack{\mathcal{QC}\!\mathit{oh}}
\def\Cohstack{\mathcal{C}\!\mathit{oh}}
\def\Spacesstack{\mathcal{S}\!\mathit{paces}}
\def\Quotfunctor{\mathrm{Quot}}
\def\Hilbfunctor{\mathrm{Hilb}}
\def\Curvesstack{\mathcal{C}\!\mathit{urves}}
\def\Polarizedstack{\mathcal{P}\!\mathit{olarized}}
\def\Complexesstack{\mathcal{C}\!\mathit{omplexes}}
% \Pic is the operator that assigns to X its picard group, usage \Pic(X)
% \Picardstack_{X/B} denotes the Picard stack of X over B
% \Picardfunctor_{X/B} denotes the Picard functor of X over B
\def\Pic{\mathop{\mathrm{Pic}}\nolimits}
\def\Picardstack{\mathcal{P}\!\mathit{ic}}
\def\Picardfunctor{\mathrm{Pic}}
\def\Deformationcategory{\mathcal{D}\!\mathit{ef}}
\usepackage{cleveref}
\crefname{originallemma}{original lemma}{original lemmas}
\Crefname{originallemma}{Original lemma}{Original lemmas}
\makeatletter
\newlabel{lemma-helper-polynomial}{{051N}{source}{}{sourcepointer-1}{}}
\newlabel{lemma-helper-polynomial@cref}{{[originallemma][0][]051N}{[1][0][]source}{}{}{}}
\makeatother
\title{A refined Noether normalization controlling the contraction of a given prime}
\author{The Stacks Project Authors}
\date{stacks-project:algebra.tex:local-snapshot-20261012; GNU FDL 1.2 or later}
\begin{document}
\maketitle
\noindent\textit{Editorial note.} All delivered mathematical text is the author's own, copied from the pinned Stacks Project chapter file algebra.tex; no statement or proof step was rewritten, reordered or removed.\par\medskip
\noindent\textit{Editorial note.} Preamble provenance, as required for this source: the frozen snapshot stages the chapter's own preamble file beside it (preamble.tex, 217 lines, 8116 bytes, the same shared Stacks-wide preamble the published chow and etale-cohomology cards use); the first line of the pinned chapter file is \textbackslash{}input\{preamble\}, so no preamble text lives inside the chapter itself. The card takes that file's lines 8-18 (the verbatim package and the three \textbackslash{}comment-based environments reference/slogan/history at lines 9-12, section 13.2, together with the Xy-pic setup) and its lines 154 to the runtime end of the file (the native theorem declarations and mathematical macros). Not one character was changed, no package was added, no macro was defined, and the excerpt serves only as preamble, never as a span coordinate.\par\medskip
\noindent\textit{Editorial note.} The same-chapter pointers are located at their \textbackslash{}label lines inside the pinned chapter file; the cross-chapter pointers are located in the project's own tag list staged in the snapshot (sources/stacks-chow/tags.txt), at the verbatim lines carrying their canonical tags.\par\medskip
\noindent\textit{Editorial note.} The counted proof refers to same-chapter statements by label; they are recorded as navigation stubs pointed at their \textbackslash{}label lines, and the arguments stay inside the counted span.\par\medskip
\noindent\textit{Editorial note.} The counted spans contain no \textbackslash{}cite, so the card carries no bibliography.\par\medskip
\noindent\textit{Allowed standard prerequisites (editorial).}\begin{itemize}
\item Induction on the number of variables. Used in the counted proof (lines 28300-28310). Reduces the statement to fewer variables after a change of coordinates.
\item Change of coordinates putting a nonzero element in a good shape. Used in the counted proof (lines 28302-28320). Makes the chosen element monic in one variable so that the quotient becomes finite.
\end{itemize}\medskip
\ifdefined\hypertarget\hypertarget{sourcepointer-1}{}\fi
\noindent\textit{Source reference.} Original lemma 051N: sources/stacks-algebra/pinned-original-source.tex, source lines 28180--28180. Navigation pointer for \textbackslash{}ref\{lemma-helper-polynomial\} used inside the counted span (line 28312); the argument that consumes it is retained in the span.\par\medskip
\begin{lemma}
\label{lemma-refined-Noether-normalization}
Let $k$ be a field. Let $\mathfrak q \subset k[x_1, \ldots, x_n]$
be a prime ideal. Set $r = \text{trdeg}_k\ \kappa(\mathfrak q)$.
Then there exists a finite ring map
$\varphi : k[y_1, \ldots, y_n] \to k[x_1, \ldots, x_n]$ such
that $\varphi^{-1}(\mathfrak q) = (y_{r + 1}, \ldots, y_n)$.
\end{lemma}

\begin{proof}
By induction on $n$. The case $n = 0$ is clear. Assume $n > 0$.
If $r = n$, then $\mathfrak q = (0)$ and the result is clear.
Choose a nonzero $f \in \mathfrak q$. Of course $f$ is nonconstant.
After applying an automorphism of the form
$$
k[x_1, \ldots, x_n] \longrightarrow k[x_1, \ldots, x_n],
\quad
x_n \mapsto x_n,
\quad
x_i \mapsto x_i + x_n^{e_i}\ (i < n)
$$
we may assume that $f$ is monic in $x_n$ over $k[x_1, \ldots, x_n]$, see
Lemma \ref{lemma-helper-polynomial}. Hence the ring map
$$
k[y_1, \ldots, y_n] \longrightarrow k[x_1, \ldots, x_n],
\quad
y_n \mapsto f,
\quad
y_i \mapsto x_i\ (i < n)
$$
is finite. Moreover $y_n \in \mathfrak q \cap k[y_1, \ldots, y_n]$ by
construction. Thus
$\mathfrak q \cap k[y_1, \ldots, y_n] = \mathfrak pk[y_1, \ldots, y_n] + (y_n)$
where $\mathfrak p \subset k[y_1, \ldots, y_{n - 1}]$ is a prime ideal.
Note that $\kappa(\mathfrak p) \subset \kappa(\mathfrak q)$ is finite, and
hence $r = \text{trdeg}_k\ \kappa(\mathfrak p)$.
Apply the induction hypothesis to the pair
$(k[y_1, \ldots, y_{n - 1}], \mathfrak p)$ and we obtain a finite ring map
$k[z_1, \ldots, z_{n - 1}] \to k[y_1, \ldots, y_{n - 1}]$ such that
$\mathfrak p \cap k[z_1, \ldots, z_{n - 1}] = (z_{r + 1}, \ldots, z_{n - 1})$.
We extend the ring map
$k[z_1, \ldots, z_{n - 1}] \to k[y_1, \ldots, y_{n - 1}]$
to a ring map
$k[z_1, \ldots, z_n] \to k[y_1, \ldots, y_n]$
by mapping $z_n$ to $y_n$.
The composition of the ring maps
$$
k[z_1, \ldots, z_n] \to k[y_1, \ldots, y_n] \to k[x_1, \ldots, x_n]
$$
solves the problem.
\end{proof}

\end{document}
